\documentclass[11pt]{article}
\usepackage[margin=1in]{geometry}
\usepackage[T1]{fontenc}
\usepackage{lmodern}
\usepackage{amsmath,amssymb,amsthm,mathtools}
\usepackage[colorlinks=true,linkcolor=blue,urlcolor=blue]{hyperref}
\urlstyle{same}
\emergencystretch=2em
\newtheorem{theorem}{Theorem}
\newtheorem{lemma}{Lemma}
\newtheorem{proposition}{Proposition}
\newtheorem{corollary}{Corollary}
\numberwithin{equation}{section}
\newcommand{\R}{\mathbb R}
\newcommand{\C}{\mathbb C}
\newcommand{\HH}{\mathbb H}
\newcommand{\dd}{\,d}
\newcommand{\Rea}{\operatorname{Re}}
\DeclareMathOperator{\cof}{cof}
\DeclareMathOperator{\tr}{tr}
\title{Sharp logarithmic estimates for Kakeya tubes}
\author{Working manuscript}
\date{November 12, 2026}
\begin{document}
\maketitle

\begin{abstract}
We give a detailed argument for a sharp logarithmic overlap estimate
for direction-separated tubes in Euclidean space. The estimate allows
arbitrary nonnegative coefficients and is equivalent to the critical
Kakeya maximal inequality with logarithmic power $(n-1)/n$.
The proof associates compact conserved positive matrix fields to
tubes. A Dirichlet determinant energy is superharmonic in three
auxiliary parameters. A mixed-determinant trace inequality gives the
weighted input bound, while a logarithmic test function gives a
matching improvement at the output. Smooth Dirichlet solvability is
taken from the theorem of Ma and Zhu on flat quaternionic quotients.
The remaining calculus, packet identities, estimates, and maximal
duality are proved below.
\end{abstract}

\section{The tube estimate}

For $a\in\R^n$, $\omega\in S^{n-1}$, and $0<\delta\le1/2$, let
$T_\delta(\omega,a)$ denote the open $\delta$-neighborhood of
$a+[-1/2,1/2]\omega$. Its volume is comparable to $\delta^{n-1}$.
Write $1_T$ for the indicator of a measurable set $T$.

\begin{theorem}[Weighted tube estimate]\label{thm:tubes}
Fix $n\ge3$ and set $q=n/(n-1)$. There is a constant $C_n$ such that,
for every $0<\delta\le1/2$, every finite $\delta$-separated set
$\{\omega_j\}\subset S^{n-1}$, every choice of centers
$a_j\in\R^n$, and every choice of coefficients $\lambda_j\ge0$,
\begin{equation}\label{eq:maintubes}
 \int_{\R^n}\left(\sum_j\lambda_j
              1_{T_\delta(\omega_j,a_j)}(X)\right)^q\dd X
 \le C_n\log(2/\delta)\,\delta^{n-1}\sum_j\lambda_j^q.
\end{equation}
The power one on the logarithm is optimal.
\end{theorem}

The analytic input used to prove this theorem is precisely the smooth
Dirichlet theorem stated in Section~\ref{sec:dirichlet}.
Its application to the present domains is proved there.
This manuscript is a research deduction and has not been independently
verified. All coefficients and tube positions in
Theorem~\ref{thm:tubes} are unrestricted apart from the stated
direction separation.

Define the Kakeya maximal operator by
\[
 K_\delta f(\omega)=\sup_{a\in\R^n}
 \frac{1}{|T_\delta(\omega,a)|}
                   \int_{T_\delta(\omega,a)}|f(X)|\dd X.
\]

\begin{corollary}\label{cor:maximal}
For $n\ge3$ and $0<\delta\le1/2$,
\begin{equation}\label{eq:maximal}
 \|K_\delta f\|_{L^n(S^{n-1})}
 \le C_n[\log(2/\delta)]^{(n-1)/n}\|f\|_{L^n(\R^n)}.
\end{equation}
This logarithmic power is optimal. Conversely,
\eqref{eq:maximal} implies \eqref{eq:maintubes}, with a change in
the dimensional constant.
\end{corollary}

For the unweighted family in Theorem~\ref{thm:tubes}, H\"older's
inequality gives
\[
 \left|\bigcup_jT_\delta(\omega_j,a_j)\right|
 \ge c_n\,\frac{\#\{j\}\delta^{n-1}}
                  {[\log(2/\delta)]^{n-1}}.
\]
In particular, every Kakeya set $K\subset\R^n$ satisfies
\begin{equation}\label{eq:neighborhood}
 |N_\delta(K)|\ge c_n[\log(2/\delta)]^{-(n-1)}
\end{equation}
for $0<\delta\le1/2$, where $N_\delta(K)$ is its open
$\delta$-neighborhood.

Here is the normalized estimate from which the theorem will follow.
Put $m=n-1$ and $q=(m+1)/m$. Fix $B_0>0$ and a sufficiently small
dimensional constant $\gamma>0$. Consider a finite family
\[
 T=\{(y,z)\in\R^m\times[0,1]:|y-b_T-zv_T|\le\delta\},
 \qquad |b_T|\le B_0,\quad |v_T|\le\gamma,
\]
with $\delta$-separated slopes $v_T$.
For $\lambda_T\ge0$ set
\[
 M(y,z)=\sum_T\lambda_T1_T(y,z),\qquad
 E=\delta^m\sum_T\lambda_T^q,\qquad L=\log(2/\delta).
\]

\begin{proposition}[Normalized tube estimate]\label{thm:main}
For sufficiently small $\delta$, depending on $m$ and $B_0$,
\begin{equation}\label{eq:overlap}
 \int_0^1\int_{\R^m}M(y,z)^q\dd y\dd z\le C_{m,B_0}EL.
\end{equation}
\end{proposition}

The proof has three quantitative steps. With $r=\delta^2$ and
$h=\delta^4$, the packet construction gives a nonnegative
superharmonic function $F(t,Y)$ for $t\in\R$ and
$Y\in\C\setminus\{0\}$ in a fixed bounded parameter region.
The weighted trace estimate gives $F(0,Y)\lesssim EL$ on $|Y|=1$.
Scalar propagation costs one logarithm. The output test gives
$F(1+rz,Y)\gtrsim r^{-1}L\int M(y,z)^q\dd y$ on $|Y|=h$.
Integration in $t$ and cancellation of one $L$ prove
Proposition~\ref{thm:main}.

The weighted input estimate uses repeated integration by parts for
mixed determinants. A related complex inequality is recorded in
\cite[proof of Lemma 5.1]{KN}; the version needed here is proved in
Lemma~\ref{lem:weightedenergy}. No decomposition into coefficient
ranges is needed.

Acknowledgments and AI usage. This manuscript was developed with an
OpenAI model. The conserved-current and determinant-energy strategy
comes from \cite{OpenAI}; its compact packets and quadratic
specialization were studied in \cite{Draft}. The logarithmic output
test and the preceding unweighted argument appear in
\cite{Addendum,AllDimensions}. The smooth Dirichlet theorem is due to
Ma and Zhu \cite{MZ}. The present manuscript supplies the weighted
input estimate and gives the full tube and maximal-function deduction.

\section{The three-coordinate Hessian and its cofactor}

Use coordinates $\zeta=(x,P)\in\R^m\times\C_{\mathrm{k}}^m$, where
$\C_{\mathrm{k}}=\R+\mathrm{k}\R$ and
$P_j=p_j+\mathrm{k} s_j$.
The three real coordinates of index $j$ are $(x_j,p_j,s_j)$.
Set $(e_0,e_1,e_2)=(1,\mathrm{i},\mathrm{j})$, where
$\mathrm{i}\mathrm{j}=\mathrm{k}$ and
$\mathrm{i}^2=\mathrm{j}^2=\mathrm{k}^2=-1$.
For a real smooth function $v$ define
\[
 G_jv=\sum_{\alpha=0}^2 e_\alpha v_{j\alpha},\qquad
 D_{jk}v=\sum_{\alpha,\beta=0}^2e_\alpha\bar e_\beta
                                      v_{j\alpha,k\beta}.
\]
Then $Dv$ is Hermitian and $\tr Dv=\Delta v$.
We call $v$ admissible if $Dv\ge0$.
The determinant is the Moore determinant. For Hermitian matrices put
$A:B=\Rea\tr(AB)$, and write $C_v=\cof Dv$.
For a matrix $A$ its real representation has entries
$A^{\R}_{j\alpha,k\beta}=\Rea(\bar e_\alpha A_{jk}e_\beta)$.
Thus $A:Dv=A^{\R}:D^2v$ and positivity is preserved.

We use the elementary Moore determinant identities
\[
 C_v:Dv=m\det Dv,\qquad
 \cof H=(\det H)H^{-1}\quad(H>0),\qquad
 D\phi(v)=\phi'(v)Dv+\phi''(v)Gv(Gv)^*.
\]
Here the cofactor is the polynomial matrix characterized by
$\left.\frac{d}{ds}\right|_{s=0}\det(H+sA)=(\cof H):A$.
It is positive semidefinite when $H\ge0$.
In particular $D(|\zeta|^2)=6I$.

\begin{lemma}[Cofactor and Green identities]\label{lem:cofactor}
For every real smooth $v$,
\begin{equation}\label{eq:cofactor}
 \sum_{j=1}^m\sum_{\alpha=0}^2
       \bar e_\alpha\partial_{j\alpha}(C_v)_{jk}=0
       \qquad(1\le k\le m).
\end{equation}
If $g,h$ are real smooth functions vanishing on the boundary of a ball,
then
\begin{equation}\label{eq:pairing}
 \int (Gg)^*C_vGh=-\int g(C_v:Dh)\in\R,
\end{equation}
and this real pairing is symmetric in $g,h$.
The divergence identity also holds for polarized cofactor matrices.
\end{lemma}
\begin{proof}
For $m=1$ the cofactor is constant. Suppose $m\ge2$.
Polarize the cofactor, a polynomial of degree $m-1$.
Given quaternionic columns $g_1,\ldots,g_{m-1}$, put
$H(t)=\sum_{\ell=1}^{m-1}t_\ell g_\ell g_\ell^*$ with $t_\ell>0$.
Its image is the span of the $g_\ell$ and its rank is at most $m-1$.
Diagonalizing $H(t)$ shows that $\cof H(t)$ kills its image:
when the rank is $m-1$ the cofactor is supported on its kernel,
and at smaller rank it is zero.
Thus every coefficient of the polynomial $\cof H(t)$ kills each
$g_\ell$. In particular this holds for its fully polarized coefficient.
That coefficient is Hermitian, so it is also killed on the left by
each $g_\ell^*$.

Apply this observation to real exponential functions
$v_\ell(\zeta)=\exp(\xi_\ell\cdot\zeta)$.
Their Hessian symbols are $g_\ell g_\ell^*$, where
$(g_\ell)_j=\sum_\alpha e_\alpha(\xi_\ell)_{j\alpha}$.
The symbol of the divergence of the polarized cofactor is
$(\sum_\ell g_\ell)^*$ times that coefficient, and hence vanishes.
The polarized divergence is a constant-coefficient multilinear
differential expression. Its symbol is a polynomial in the
$\xi_\ell$; vanishing for all real $\xi_\ell$ proves the identity
for arbitrary smooth inputs. Taking equal inputs proves
\eqref{eq:cofactor}.

Integrate the derivative of $g$ once in the left side of
\eqref{eq:pairing}. The differentiated-cofactor term vanishes by
\eqref{eq:cofactor}, and the boundary term vanishes because $g=0$.
The remaining expression is
\[
 -\int g\sum_{j,k,\alpha,\beta}
         \bar e_\alpha(C_v)_{jk}e_\beta h_{j\alpha,k\beta}.
\]
The inner sum is real: quaternionic conjugation followed by
interchanging $(j,\alpha)$ and $(k,\beta)$ leaves it unchanged.
Its real value is $C_v:Dh$ by the definition of the real
representation. Finally, interchanging $g,h$ conjugates the original
pairing, so its real value is symmetric.
\end{proof}

The real representation of each mixed cofactor is symmetric and
divergence free. Consequently, when a real smooth function $f$ and
its first derivatives vanish on the boundary, the corresponding
linear operator can be transferred from any smooth function to $f$
by two integrations by parts, without a boundary term. We use this
observation in the weighted trace inequality.

\section{Smooth Dirichlet solutions on ordinary balls}\label{sec:dirichlet}

We use the following exact part of Ma--Zhu's theorem: on a connected
smooth relatively compact domain in $\HH^m/\Gamma$, where $\Gamma$
is any discrete additive subgroup and $m\ge2$, a smooth positive
prescribed density and smooth boundary data have a unique smooth
admissible Dirichlet solution if a smooth strictly plurisubharmonic
subsolution with those boundary data exists.

\begin{proposition}\label{prop:dirichlet}
For a Euclidean ball $B\subset\R^{3m}$ and $\mu>0$ smooth on
$\overline B$, there is a unique smooth solution
\[
 Du>0,\qquad\det Du=\mu\quad\hbox{in }B,
 \qquad u|_{\partial B}=0.
\]
It is nonpositive and depends smoothly on smooth finite-dimensional
families of strictly positive densities. If $v$ is smooth and
admissible on a smooth subdomain, $\det Dv\ge\det Du$ there, and
$v\le u$ on its boundary, then $v\le u$ throughout that subdomain.
\end{proposition}
\begin{proof}
Adjoin fourth coordinates $w\in\R^m/\mathbb Z^m$ and consider
\[
 \widetilde B=B\times(\R^m/\mathbb Z^m)
 \Subset\HH^m/(\mathrm{k}\mathbb Z^m).
\]
The ball coordinates occupy the $1,\mathrm{i},\mathrm{j}$ slots.
This domain is connected, relatively compact, and has smooth nonempty
boundary. If $B=B_R(\zeta_0)$, the globally defined function
$r(\zeta,w)=|\zeta-\zeta_0|^2-R^2$ has full quaternionic Hessian
$H(r)=6I$. Consequently $Cr$ is a strict plurisubharmonic subsolution
with zero boundary data and determinant $(6C)^m\ge\sup\mu$ for
sufficiently large $C$.

Extend $\mu$ constantly in $w$ and apply \cite[Theorem 1.1]{MZ}.
Uniqueness makes the solution invariant under all translations in $w$.
It descends to $B$, and the full Hessian there becomes exactly $D$.
Its Laplacian is positive, so its zero boundary values imply $u\le0$.

For comparison, the difference of two log determinants is an average
of positive inverse Hessians contracted with the Hessian difference.
This is a real elliptic operator, and the maximum principle gives
$v\le u$ when $\det Dv\ge\det Du$ and $v\le u$ on the boundary.
The solution has positive Hessian on the closed ball: its Hessian
is nonnegative at the boundary by continuity, and its determinant
equals the positive boundary value of $\mu$ there. A semidefinite
comparison function can be replaced by
$v+\eta(|\zeta-\zeta_0|^2-R^2)$ on a containing ball, with the
constant chosen so that the added quadratic is nonpositive on the
subdomain; then let $\eta\downarrow0$.
At any fixed positive source the linearization
$(Du)^{-1}:Dh$ is uniformly elliptic on the closed ball. Its zero-boundary
Dirichlet map is an isomorphism in H\"older spaces. The implicit
function theorem and elliptic bootstrapping give smooth parameter
dependence. No uniform estimate as $\inf\mu\downarrow0$ is needed.
\end{proof}

\section{The scalar determinant energy}

For any smooth nonnegative source on a ball define
\[
 \mathcal C_B=\{v\in C^\infty(\overline B):v|_{\partial B}=0,\ Dv\ge0\}.
\]
Set
\[
 J_\mu(v)=\int_B(-v)\mu-\frac1{m+1}\int_B(-v)\det Dv,
 \qquad F_B(\mu)=\sup_{v\in\mathcal C_B}J_\mu(v).
\]
Every admissible test is nonpositive because its ordinary Laplacian
is nonnegative and its boundary value is zero.
For smooth variations $g$ with zero boundary value, the Green
identity and $C_v:Dv=m\det Dv$ give
\[
 \left.\frac{d}{ds}\right|_{s=0}
       \int_B[-(v+sg)]\det D(v+sg)
       =-(m+1)\int_B g\det Dv.
\]
Indeed the differentiated determinant contributes
$-\int v(C_v:Dg)=-\int g(C_v:Dv)$.
It follows that
\[
 dJ_\mu(v)[g]=\int g(\det Dv-\mu),\qquad
 d^2J_\mu(v)[g,g]=-\int(Gg)^*C_vGg\le0.
\]
The admissible class is convex, so the second variation proves
concavity along its line segments. The positive-source Dirichlet
solution is a stationary point and hence a global maximizer. Thus
Proposition~\ref{prop:dirichlet} gives
\begin{equation}\label{eq:value}
 F_B(\mu)=J_\mu(u)=\frac m{m+1}\int_B(-u)\mu.
\end{equation}

The definition shows that $F_B$ is monotone and convex in the source,
and $F_B(c\mu)=c^qF_B(\mu)$ for $c\ge0$.
The source $1$ has the solution $(|\zeta-\zeta_0|^2-R^2)/6$ on $B_R$,
so $F_B(1)<\infty$. Hence for $0\le\mu,\nu\le A$ and
$a=\|\mu-\nu\|_\infty\le1$,
\begin{equation}\label{eq:continuity}
 F_B(\mu)\le F_B(\nu+a)
 \le(1-a)F_B(\nu)+aF_B(\nu+1)
 \le F_B(\nu)+a(A+1)^qF_B(1).
\end{equation}
Interchanging the sources gives local Lipschitz continuity, including
at sources that vanish.

Let $p=(p_0,p_1,p_2)\in\R^3$ be parameters. Include this triple as
coordinate index $0$, and suppose
\[
 A=\sum_\ell\rho_\ell
       \begin{pmatrix}1\\b_\ell\end{pmatrix}
       \begin{pmatrix}1\\b_\ell\end{pmatrix}^{\!*}
   =\begin{pmatrix}\mu&\xi^*\\\xi&S\end{pmatrix},\qquad\rho_\ell\ge0,
\]
is smooth on an open parameter set times $\overline B$.
Real-column conservation means
\begin{equation}\label{eq:realconservation}
 \sum_{a=0}^m\sum_{\alpha=0}^2
 \partial_{a\alpha}\Rea(\bar e_\alpha A_{ab}e_\beta)=0
 \qquad(0\le b\le m,\ 0\le\beta\le2),
\end{equation}
where $\partial_{0\alpha}=\partial_{p_\alpha}$.
Assume also that $S$ vanishes in a spatial boundary collar locally
uniformly in $p$. Positivity then forces $\xi=0$ in that collar.

\begin{proposition}[Scalar transfer]\label{prop:transfer}
The function $F(p)=F_B(\mu(p,\cdot))$ is nonnegative, locally Lipschitz,
and distributionally superharmonic. If $v$ is a smooth admissible
barrier with $\det Dv\ge\mu$ and $v\le0$ on $\partial B$, then
\begin{equation}\label{eq:barrierupper}
 F(p)\le\frac m{m+1}\int_B(-v)\mu.
\end{equation}
For pairwise disjoint balls $B_j\subset B$,
\begin{equation}\label{eq:localization}
 F_B(\mu)\ge\sum_jF_{B_j}(\mu|_{B_j}).
\end{equation}
\end{proposition}
\begin{proof}
First suppose $\mu>0$ on the closed ball. Let $u(p,\zeta)$ be its
Dirichlet solution and write
\[
 H=Du,\quad C=\cof H=\mu H^{-1},\quad
 g_\beta=\partial_{p_\beta}u,\quad
 w=\sum_{\beta=0}^2Gg_\beta\bar e_\beta.
\]
Smooth parameter dependence justifies all derivatives below.
Each $g_\beta$ has zero boundary value.
Differentiating the equation gives
$\partial_{p_\beta}\mu=C:Dg_\beta$, and hence
\begin{equation}\label{eq:parameterpairing}
 Q:=\sum_\beta\int_B(Gg_\beta)^*CGg_\beta
   =-\sum_\beta\int_Bg_\beta\partial_{p_\beta}\mu.
\end{equation}
By Lemma~\ref{lem:cofactor}, the integrated pairings
$\int(Gg_\beta)^*CGg_\gamma$ are real and symmetric.
Expanding $w^*Cw$ and using
$e_\beta\bar e_\gamma+e_\gamma\bar e_\beta=0$ for
$\beta\ne\gamma$ therefore gives
\begin{equation}\label{eq:threepairings}
 \int_Bw^*Cw=Q.
\end{equation}

The $b=0$ columns of \eqref{eq:realconservation} say
\begin{equation}\label{eq:massconservation}
 \partial_{p_\beta}\mu+
 \sum_{j,\alpha}\partial_{j\alpha}
               \Rea(\bar e_\alpha\xi_j e_\beta)=0.
\end{equation}
Differentiating these equations in $p_\beta$, summing over $\beta$,
and using the remaining columns of \eqref{eq:realconservation}
gives
\begin{equation}\label{eq:doubledconservation}
 \Delta_p\mu
 =\sum_{j,k,\alpha,\beta}\partial_{j\alpha}\partial_{k\beta}
                                      S^\R_{j\alpha,k\beta}.
\end{equation}
The passage between the two sets of columns uses
$\Rea(\bar e_\alpha\xi_j e_\beta)
 =\Rea(\bar e_\beta\bar\xi_j e_\alpha)$.
Also, integrating \eqref{eq:massconservation} against $g_\beta$
and summing gives, with \eqref{eq:parameterpairing},
\begin{equation}\label{eq:crosspairing}
 \int_B\Rea(w^*\xi)
 =\sum_\beta\int_Bg_\beta\partial_{p_\beta}\mu=-Q.
\end{equation}

Since the first variation of $J_{\mu(p)}$ vanishes at $u(p)$,
\[
 \partial_{p_\beta}F=-\int_Bu\,\partial_{p_\beta}\mu,\qquad
 \Delta_pF=Q-\int_Bu\,\Delta_p\mu.
\]
Two spatial integrations in \eqref{eq:doubledconservation} have no
boundary terms because $S$ vanishes in a collar. They give
$\int u\,\Delta_p\mu=\int S:H$.
Combining this with \eqref{eq:threepairings} and
\eqref{eq:crosspairing}, we obtain the exact identity
\begin{align}
 \Delta_pF
 &=-\int_B\big(S:H+w^*Cw+2\Rea(w^*\xi)\big)\nonumber\\
 &=-\int_B\sum_\ell\rho_\ell
       (b_\ell+H^{-1}w)^*H(b_\ell+H^{-1}w)\le0.
       \label{eq:energyidentity}
\end{align}
The last equality uses
$S=\sum_\ell\rho_\ell b_\ell b_\ell^*$,
$\xi=\sum_\ell\rho_\ell b_\ell$, and
$\mu=\sum_\ell\rho_\ell$.

For an upper barrier, comparison gives $v\le u$, so
\eqref{eq:value} proves \eqref{eq:barrierupper}.
For localization, let $u_j$ be the zero-boundary solution on $B_j$
for the restricted source. The global solution satisfies $u\le0=u_j$
on $\partial B_j$ and has the same determinant there and inside.
Comparison gives $u\le u_j$ in $B_j$.
Since $-u\ge0$, integration over the disjoint balls and
\eqref{eq:value} give
\[
 F_B(\mu)=\frac{m}{m+1}\int_B(-u)\mu
 \ge\sum_j\frac{m}{m+1}\int_{B_j}(-u_j)\mu
 =\sum_jF_{B_j}(\mu|_{B_j}).
\]

Finally let $\mu\ge0$. Add
$\epsilon\operatorname{diag}(1,0,\ldots,0)$ to the current.
This is a constant conserved positive field; it changes the source
to $\mu+\epsilon$ and leaves $S$ unchanged. Equivalently it adds
one term with $\rho=\epsilon$ and $b=0$.
The continuity estimate \eqref{eq:continuity} gives locally uniform
convergence of the scalar energies as $\epsilon\downarrow0$.
Nonnegativity, distributional superharmonicity, and finite-ball
localization pass to this limit.
For the barrier assertion set
$v_\epsilon=v+\epsilon^{1/m}(|\zeta-\zeta_0|^2-R^2)/6$ on $B_R$.
The nonnegative mixed determinants imply
$\det Dv_\epsilon\ge\det Dv+\epsilon$.
Apply the positive-source barrier bound and let $\epsilon\downarrow0$.
\end{proof}

\section{Compact packets in arbitrary dimension}

Use parameters $(t,Y)\in\R\times\C_{\mathrm{k}}$, encoded by
$t+Y\mathrm{i}$, and spatial triples $x_j+P_j\mathrm{i}$.
Set
\[
 r=\delta^2,\quad h=\delta^4,\quad \tau=t-1,\quad s=|Y|>0,
 \quad K_0=r+s,\quad D_0=K_0^2+\tau^2.
\]
Also put $W=P/Y=\alpha+\mathrm{k}\beta$ and $\varepsilon=Y/s$.
For a tube set $a=rb_T$, $v=v_T$, and define
\[
 \Phi_T=\min_{w\in\R^m}\left(
 r|w-v|^2+\frac{|x-a-\tau w|^2}{K_0}
                 +s|\alpha-w|^2+s|\beta|^2\right).
\]
Let $c$ be the minimizer and $d=(x-a-\tau c)/K_0$. Equivalently,
\[
 K_0c-\tau d=rv+s\alpha,\qquad \tau c+K_0d=x-a.
\]
Define a quaternionic column $B$ of length $m+1$ by
\[
 B_0=1,\qquad B_j=\tfrac12(W_j+c_j-\varepsilon\mathrm{i} d_j),
 \qquad f_0=(s^2D_0)^{-m/2}.
\]

\begin{lemma}\label{lem:packets}
With $G$ now including the parameter triple,
\begin{equation}\label{eq:packetidentities}
 \sum_{j=0}^m\sum_{\alpha=0}^2
       \bar e_\alpha\partial_{j\alpha}(f_0B_j\bar B_k)=0
       \quad(0\le k\le m),\qquad B^*G\Phi_T=0.
\end{equation}
\end{lemma}
\begin{proof}
It suffices to calculate at $Y=s>0$.
Indeed, for $\eta=\exp(\mathrm{k}\theta)$ the rotation
$(Y,P)\mapsto(\eta Y,\eta P)$ preserves $\Phi_T$, $s$, $W$, and $f_0$.
With $u=\exp(\mathrm{k}\theta/2)$, it sends each encoded triple
$x+P\mathrm{i}$ to $u(x+P\mathrm{i})\bar u$, and sends each $B_j$
to $uB_j\bar u$. The gradient and divergence expressions in
\eqref{eq:packetidentities} transform by the same conjugation.
Thus their vanishing is preserved under this rotation.
For each spatial coordinate, regard $c,d$ first as functions of
$(\tau,s,x,\alpha)$, with the other quantities fixed.
Differentiation in these coordinates gives
\[
 c_\tau+d_s+\alpha c_x=0,\quad
 d_\tau-c_s+\alpha d_x=0,\quad
 c_\alpha=sd_x,\quad d_\alpha=-sc_x,
 \qquad c_x=\tau/D_0,\quad d_x=K_0/D_0.
\]
These identities follow directly from
$c+\mathrm i d=(rv+s\alpha+\mathrm i(x-a))/(K_0+\mathrm i\tau)$,
where $\mathrm i$ in this formula is an ordinary complex unit.
For an actual derivative in $\Rea Y$ with $P$ fixed at $Y=s>0$,
the chain rule replaces $\partial_s$ by
$\partial_s-(\alpha/s)\partial_\alpha-(\beta/s)\partial_\beta$;
$c,d$ themselves are independent of $\beta$ in these coordinates.
At $Y=s$, also
$\alpha_{\operatorname{Im}_{\mathrm{k}}Y}=\beta/s$,
$\beta_{\operatorname{Im}_{\mathrm{k}}Y}=-\alpha/s$, and
$\varepsilon_{\operatorname{Im}_{\mathrm{k}}Y}=\mathrm{k}/s$.

For each $j\ge1$,
\[
 \sum_{\alpha=0}^2\bar e_\alpha\partial_{j\alpha}B_j
   =\tau/D_0-\mathrm{i}(K_0/D_0+1/s).
\]
These $m$ terms cancel the time derivatives of $f_0$, since
$(\log f_0)_\tau=-m\tau/D_0$ and
$(\log f_0)_s=-m/s-mK_0/D_0$.
This proves the first-column identity.

For column $k\ge1$, the product rule leaves the acceleration
$\sum_{j,\alpha}\bar e_\alpha B_j\partial_{j\alpha}\bar B_k$.
Its spatial terms with $j\ne k$ vanish.
Its time contribution is
\[
 \tfrac12\big(d_k/s-\mathrm{j}\beta_k(d_k)_{x_k}
                         -\mathrm{k}\beta_k(c_k)_{x_k}\big).
\]
Writing $B=(\alpha_k+c_k-\mathrm{i} d_k+\mathrm{k}\beta_k)/2$,
the remaining spatial term is
\[
 \tfrac12\big[(c_k)_{x_k}(B+\mathrm{i} B\mathrm{i})
 +(d_k)_{x_k}(B\mathrm{i}-\mathrm{i} B)
 +(\mathrm{j}B\mathrm{k}-\mathrm{i} B)/s\big],
\]
which is its negative because the three expressions in parentheses
are $\mathrm{k}\beta_k$, $\mathrm{j}\beta_k$, and $-d_k$.

Finally differentiation of the minimum gives
\[
 G_j\Phi_T=2[d_j+\mathrm{i}(\alpha_j-c_j)+\mathrm{j}\beta_j],
\]
\[
 G_0\Phi_T=-2c\cdot d+
 \mathrm{i}(|c|^2-|d|^2-|\alpha|^2-|\beta|^2)-2\mathrm{j} c\cdot\beta.
\]
Direct multiplication yields
\[
 \bar B_jG_j\Phi_T=2c_jd_j+
 \mathrm{i}(\alpha_j^2+\beta_j^2-c_j^2+d_j^2)+2\mathrm{j} c_j\beta_j.
\]
Summing proves tangency. All calculations are independent of $m$
apart from the first-column sum.
\end{proof}

Choose a fixed nonnegative smooth compactly supported function $\chi$
on $[0,\infty)$ equal to one on a sufficiently large fixed interval,
and set $Z_\chi=\int_{\R^{3m}}\chi(|u|^2)\dd u$. Define
\begin{equation}\label{eq:normalizedpacket}
 A_T=\frac{\delta^m r^{m/2}}{Z_\chi h^{3m/2}}
       f_0\chi(\Phi_T/h)BB^*,\qquad \mu_T=(A_T)_{00}.
\end{equation}
The cutoff preserves the strong conservation identity:
the additional differentiated-cutoff term in column $k$ is a real
multiple of
$\overline{B^*G\Phi_T}\,\bar B_k$, which vanishes.
Taking real parts after multiplication by each $e_\beta$ gives
\eqref{eq:realconservation}. Completing the square gives
\[
 \Phi_T=\frac{rs}{K_0}|\alpha-v|^2
 +\frac{K_0}{D_0}\left|x-a-\frac{\tau(rv+s\alpha)}{K_0}\right|^2
 +s|\beta|^2.
\]
Since $\dd P=s^{2m}\dd\alpha\dd\beta$, the change of variables yields
\[
 \int\chi(\Phi_T/h)\dd x\dd P
   =Z_\chi h^{3m/2}s^mD_0^{m/2}r^{-m/2},\qquad
 \int\mu_T=\delta^m.
\]
On packet support,
\[
 |c-v|\lesssim\sqrt{h/r},\quad |d|\lesssim\sqrt{h/K_0},\quad
 |\alpha-c|+|\beta|\lesssim\sqrt{h/s}.
\]
Thus for bounded $t,s$, both $|x-a-\tau v|$ and $|P-Yv|$ are
$O(\delta)$. Work on $-3<t<5$, $0<s<4$. For sufficiently small
$\gamma$, then small $\delta$ depending on $B_0$, every packet lies
in the spatial ball $B_{1/2}\subset\R^{3m}$ throughout this domain.
The weighted sum $A=\sum_T\lambda_TA_T$ therefore satisfies
Proposition~\ref{prop:transfer} on $B_1$. Its source is
$\mu=\sum_T\lambda_T\mu_T$ and has mass
$\delta^m\sum_T\lambda_T$ at every parameter point. This mass is
different from $E$; the following estimate gives the required $E$
dependence directly.

\section{A weighted energy estimate}

For Hermitian matrices write
$\mathcal D(A_1,\ldots,A_m)$ for the symmetric multilinear
polarization of the Moore determinant, normalized by
$\mathcal D(A,\ldots,A)=\det A$.
Repeated arguments are indicated by brackets, for example
$\mathcal D((Dv)[j],(D\phi)[m-j])$.
Mixed determinants of positive semidefinite matrices are nonnegative
\cite[Lemma 2.38]{MZ}. This also follows by decomposing each input
into nonnegative rank-one matrices and using
$\mathcal D(g_1g_1^*,\ldots,g_mg_m^*)=\det(GG^*)/m!$,
where $G$ has columns $g_1,\ldots,g_m$.
Polarization of \eqref{eq:cofactor} shows that the mixed cofactor
matrices are divergence free as well.

\begin{lemma}[Trace and weighted energy]\label{lem:weightedenergy}
Let $\phi\in C^\infty(\overline B)$ be admissible and
$-H\le\phi\le0$, where $H\ge0$.
For every $v\in\mathcal C_B$,
\begin{equation}\label{eq:trace}
 \int_B(-v)^{m+1}\det D\phi
 \le (m+1)!\,H^m\int_B(-v)\det Dv.
\end{equation}
Consequently, if a nonnegative source satisfies
$\mu\le b\det D\phi$ for a nonnegative measurable $b$, then
\begin{equation}\label{eq:weightedenergy}
 F_B(\mu)\le C_mH\int_B b^q\det D\phi.
\end{equation}
No boundary condition on $\phi$ is needed.
\end{lemma}
\begin{proof}
Put $U=-v\ge0$ and, for $0\le j\le m$, define
\[
 I_j=\int_B U^{j+1}
           \mathcal D((Dv)[m-j],(D\phi)[j]).
\]
Thus $I_0=\int U\det Dv$ and $I_m$ is the left side of
\eqref{eq:trace}. For $j\ge1$, integrate twice in one $D\phi$
factor. The mixed cofactor identity gives
\[
 I_j=\int_B\phi\,
 \mathcal D(D(U^{j+1}),(Dv)[m-j],(D\phi)[j-1]).
\]
There are no boundary terms: $U^{j+1}=0$ and
$\nabla(U^{j+1})=0$ on $\partial B$ because $j+1\ge2$.
The remaining factors may be nonzero there.
The chain rule gives
\[
 D(U^{j+1})=-(j+1)U^jDv
             +j(j+1)U^{j-1}Gv(Gv)^*.
\]
The mixed determinant involving the last, positive semidefinite,
term is nonnegative. Since $\phi\le0$, its contribution to $I_j$
is nonpositive. The other contribution is bounded by
\[
 I_j\le (j+1)H I_{j-1}.
\]
Iteration proves \eqref{eq:trace}. This argument uses only smooth
functions and remains valid when one of the Hessians is singular.

Let $V=\int b^q\det D\phi$ and
$\mathcal A(v)=\int(-v)\det Dv$. H\"older's inequality and
\eqref{eq:trace} give
\begin{align*}
 \int(-v)\mu
 &\le V^{m/(m+1)}
      \left(\int(-v)^{m+1}\det D\phi\right)^{1/(m+1)}\\
 &\le [(m+1)!]^{1/(m+1)}
          H^{m/(m+1)}V^{m/(m+1)}
          \mathcal A(v)^{1/(m+1)}.
\end{align*}
Subtract $\mathcal A(v)/(m+1)$ and maximize over
$\mathcal A(v)\ge0$. The resulting upper bound is
\[
 \frac{m}{m+1}[(m+1)!]^{1/m}HV.
\]
Taking the supremum over $v\in\mathcal C_B$ proves
\eqref{eq:weightedenergy}. The cases $H=0$ or $V=0$ follow directly;
if $V=\infty$ there is nothing to prove.
\end{proof}

\section{The weighted input bound and parameter propagation}

At $t=0$, $|Y|=1$, packet support implies
\[
 |\alpha-v_T|\le C\delta,\qquad
 |x+\alpha|+|\beta|\le C\sqrt h,
\]
where $C$ may depend on the fixed intercept bound $B_0$.
Indeed $|w-v_T|\lesssim\delta$ for a minimizer,
$|\alpha-w|+|\beta|\lesssim\sqrt h$, and, since $\tau=-1$,
$|x-rb_T+w|\lesssim\sqrt h$.
Define the directional weight function
\[
 b_\delta(\alpha)=\sum_T\lambda_T
                        1_{\{|\alpha-v_T|\le C\delta\}}.
\]
At these parameter values the scalar prefactor in
\eqref{eq:normalizedpacket} is $O(h^{-m})$.
Consequently, with
$\rho=|x+\alpha|^2+2|\beta|^2$,
\begin{equation}\label{eq:inputdomination}
 \mu\le C h^{-m}b_\delta(\alpha)1_{\{\rho\le Ch\}}.
\end{equation}

At $|Y|=1$ one has $D\rho=8I$, $|G\rho|^2=8\rho$, and
$\rho\le2$ on $B_1$. Thus the same logarithmic barrier as in the
unweighted argument,
\[
 \phi_h=C_0\log\frac{h+\rho}{h+2},
\]
is admissible, nonpositive on $B_1$, and satisfies
\[
 H_h:=\|\phi_h\|_\infty
       =C_0\log\frac{h+2}{h}\lesssim L,\qquad
 \det D\phi_h=(8C_0)^m\frac{h}{(h+\rho)^{m+1}}.
\]
Choose the fixed $C_0$ large enough that
\eqref{eq:inputdomination} implies
$\mu\le b_\delta(\alpha)\det D\phi_h$.
This choice is independent of every coefficient $\lambda_T$.

The determinant measure has uniformly bounded mass in the normal
variables. At fixed $\alpha$, translate $u=x+\alpha$ and rescale
$\beta$ by $\sqrt2$ to obtain the exact whole-space identity
\[
 \int_{\R^m\times\R^m}
 \frac{h}{(h+|x+\alpha|^2+2|\beta|^2)^{m+1}}
       \dd x\dd\beta
 =2^{-m/2}\frac{\pi^m}{m!}.
\]
Extending the positive integral from $B_1$ to these normal spaces
therefore gives
\[
 \int_{B_1}b_\delta(\alpha)^q\det D\phi_h
 \le C_{m,B_0}\int_{\R^m}b_\delta(\alpha)^q\dd\alpha
 \le C_{m,B_0}\delta^m\sum_T\lambda_T^q=C_{m,B_0}E.
\]
The last inequality follows because $\delta$-separation bounds the
number of balls $\{|\alpha-v_T|\le C\delta\}$ through each point by
a dimensional constant depending also on the fixed $C$.
Lemma~\ref{lem:weightedenergy} now proves
\begin{equation}\label{eq:input}
 F(0,Y)\le C E L\qquad(|Y|=1).
\end{equation}
This is the only additional estimate needed for arbitrary coefficients.

\begin{lemma}[Scalar propagation]\label{lem:propagation}
There are absolute constants $h_0>0$ and $C$ such that, for
$0<h<h_0$ and any nonnegative continuous distributionally
superharmonic function $G$ on
$(-2,4)\times\{Y\in\C:h/2<|Y|<3\}$,
\begin{equation}\label{eq:propagation}
 \int_1^2\operatorname{avg}_{|Y|=h}G(t,Y)\dd t
 \le C\log(2/h)\operatorname{avg}_{|Y|=1}G(0,Y).
\end{equation}
The circle averages are with respect to normalized angular measure.
\end{lemma}
\begin{proof}
We first treat smooth $G$. Circular averaging produces a function
$H(t,s)$ satisfying
$H_{tt}+H_{ss}+s^{-1}H_s\le0$.
Set $\nu=\pi/3$ and
$a(s)=\int_0^3\sin(\nu t)H(t,s)\dd t$.
Twice integrating in $t$ gives
\[
 a''+s^{-1}a'-\nu^2a
 \le-\nu[H(0,s)+H(3,s)]\le0.
\]
Fix $R_0=1/4$ and write $\ell(s)=\log(R_0/s)$.
The function
\[
 b(s)=\left(1+\frac{\nu^2s^2}{4}\right)\ell(s)
                       +\frac{\nu^2}{4}(s^2-R_0^2)
\]
vanishes at $R_0$ and satisfies, for $0<s<R_0$,
\[
 (1-\nu^2R_0^2/4)\ell(s)\le b(s)\le\ell(s),\qquad
 b''+s^{-1}b'-\nu^2b
 =\frac{\nu^4}{4}
       [R_0^2-s^2-s^2\ell(s)]\ge0.
\]
For the first bounds use
$0\le1-e^{-2u}(u+1)\le u$ for $u\ge0$.
The last sign follows from $e^{-2u}(1+u)\le1$.
Comparison for the operator
$d^2/ds^2+s^{-1}d/ds-\nu^2$ on $[h,R_0]$ gives
$a(s)\ge a(h)b(s)/b(h)$.
On the fixed interval $[R_0/3,R_0/2]$, $b$ is bounded below,
so
\[
 a(h)\le C\log(2/h)\int_{R_0/3}^{R_0/2}a(s)\dd s.
\]

We record why the last integral is controlled by $H(0,1)$.
The circularly averaged function
$\overline G(t,Y)=H(t,|Y|)$ is nonnegative and superharmonic.
The fixed compact set
$[0,3]\times\{R_0/3\le|Y|\le R_0/2\}$ and the point $(0,1)$
can be joined and covered by finitely many balls whose concentric
quadruples lie in a fixed connected region away from $Y=0$.
The mean inequality first controls the integral on a ball centered
at $(0,1)$ by a constant times $H(0,1)$.
For consecutive equal-radius balls with nearby centers, select a
point in the first where the function does not exceed its average.
A ball of three times the radius centered at that point contains
the second ball and remains in the region. The mean inequality and
nonnegativity control the integral on the second by a constant times
the integral on the first. Iterating over the finite covering
controls the integral on the compact set by $CH(0,1)$.
All radii and the number of balls are independent of $h$.
Since $s$ stays in a fixed positive interval there, cylindrical
volume is comparable to $\dd t\dd s\dd\theta$.
Consequently
$\int_{R_0/3}^{R_0/2}a(s)\dd s\le CH(0,1)$.

On $1\le t\le2$, $\sin(\nu t)\ge\sqrt3/2$, so the bound for $a(h)$
proves \eqref{eq:propagation}.
For continuous distributionally superharmonic $G$, convolve with a
nonnegative smooth approximate identity on a slightly smaller
region still containing the points and compact sets just used.
The convolutions are nonnegative and superharmonic and converge
uniformly on those compact sets. Letting the convolution radius
tend to zero proves the assertion.
\end{proof}

Applying Lemma~\ref{lem:propagation} to the scalar energy and using
\eqref{eq:input} yields
\begin{equation}\label{eq:totalupper}
 \int_1^2\operatorname{avg}_{|Y|=h}F(t,Y)\dd t\le C E L^2.
\end{equation}

\section{The logarithmic output test}

\begin{lemma}\label{lem:thin}
Let $B=B_R((x_0,0))\subset\R^m\times\C_{\mathrm{k}}^m$ and
$0<\epsilon<R/10$. If $\mu\ge a\ge0$ on
$|x-x_0|\le R/2$, $|P|\le\epsilon$, then
\begin{equation}\label{eq:thin}
 F_B(\mu)\ge c_mR^m\epsilon^{2m+2}a^q\log(R/\epsilon).
\end{equation}
\end{lemma}
\begin{proof}
Put $\eta=\epsilon/R$. On the unit ball take
\[
 v_\eta=\log\frac{b}{d},\qquad
 b=\eta^2+|P|^2,\quad d=\eta^2+1-|x|^2.
\]
It is smooth, nonpositive, and zero on the sphere. Its Hessian is
\[
 Dv_\eta=\frac4b I-\frac{gg^*}{b^2}
                    +\frac2d I+\frac{4xx^{\top}}{d^2},
 \qquad g=G(|P|^2),\quad |g|^2=4|P|^2.
\]
The first two terms together have eigenvalues $4/b$ with multiplicity
$m-1$ and $4\eta^2/b^2$ with multiplicity one. Thus $Dv_\eta>0$.
Since $b\le d$, add $2I/d$ first, then the rank-one last term, to get
\[
 \det Dv_\eta\le C_m\left(
 \frac{\eta^2}{b^{m+1}}+
 \frac1{d b^{m-1}}+\frac{|x|^2}{d^2b^{m-1}}\right).
\]
In fact the determinant before the last addition is
$(4\eta^2/b^2+2/d)(4/b+2/d)^{m-1}$, and its cofactor norm is at
most $(6/b)^{m-1}$.

For fixed $x$, write $s_0=1-|x|^2$, $t_0=|P|^2\in[0,s_0]$.
Radial integration in $P$ has measure $c_mt_0^{m-1}\dd t_0$.
For the first term use
\[
 \int_0^{s_0}\frac{t_0^{m-1}\eta^2}{(\eta^2+t_0)^{m+1}}\dd t_0
 \le C_m,
\]
and bound $-v_\eta$ by $\log((\eta^2+s_0)/\eta^2)$.
For the other two terms use
\[
 \int_0^{s_0}\frac{t_0^{m-1}}{(\eta^2+t_0)^{m-1}}
       \log\frac{\eta^2+s_0}{\eta^2+t_0}\dd t_0\le s_0.
\]
The remaining $x$ integral is bounded by
\[
 C_m\int_{|x|<1}s_0\left(\frac1{\eta^2+s_0}
                  +\frac{|x|^2}{(\eta^2+s_0)^2}\right)\dd x
 \le C_m\log(1/\eta).
\]
Near $|x|=1$, radial measure is comparable to $\dd s_0$; away from
that boundary the integral is bounded. Consequently
\[
 \int_{B_1}(-v_\eta)\det Dv_\eta\le C_m\log(1/\eta).
\]
On $|x|\le1/2$, $|P|\le\eta$, its depth is at least
$c_m\log(1/\eta)$.
The test $A_0v_\eta((x-x_0)/R,P/R)$ in the variational functional
therefore has value at least
\[
 R^m\log(R/\epsilon)
          [c_m A_0a\epsilon^{2m}-C_m A_0^{m+1}].
\]
Choose $A_0=c'_ma^{1/m}\epsilon^2$ with $c'_m>0$ small.
This proves \eqref{eq:thin}; the zero case follows from $F_B\ge0$.
\end{proof}

\section{Output localization and the final cancellation}

Fix $z\in[0,1]$ and set $t=1+rz$, $|Y|=h$.
Partition $\R^m$ into cubes $Q$ of side $\delta$ with centers $y_Q$.
Let $M_Q=\sup_{y\in Q}M(y,z)$, choose a point attaining this
value, and select the tubes present there. The maximum is attained
because a finite weighted family gives only finitely many possible
values of $M$. For the finitely many cubes with $M_Q>0$, take
disjoint spatial balls
\[
 B_Q=B_R((r y_Q,0)),\qquad R=c_0r\delta,
\]
with $c_0=1/10$. These balls are pairwise disjoint.
Since $|y_Q|\le B_0+\gamma+C_m\delta$ for the cubes in question,
they all lie in $B_{1/2}$ for sufficiently small $\delta$. On
\[
 |x-r y_Q|\le R/2,\qquad |P|\le\epsilon=h/100,
\]
the packets satisfy
\begin{equation}\label{eq:outputsource}
 \mu(t,Y,x,P)\ge c h^{-5m/2}M_Q.
\end{equation}
Indeed, for each selected tube choose $w=v_T$ in $\Phi_T$. The tube condition at the chosen
point of $Q$ gives $|x-rb_T-rzv_T|\le Cr\delta$, so its quadratic
term is at most $Cr\delta^2=Ch$. Also $|P/Y|\le1/100$,
and the other terms are at most $Ch$. Choose the fixed flat interval
of $\chi$ to contain this range. At the output,
$D_0\asymp r^2=h$, so the scalar prefactor in
\eqref{eq:normalizedpacket} is comparable to
$h^{-m}(h^3)^{-m/2}=h^{-5m/2}$.
Multiplying the lower bound for each selected packet by $\lambda_T$
and summing gives \eqref{eq:outputsource}.

The source cylinder just used lies inside $B_Q$.
Since $\epsilon/R=\delta/10$, one has
$\epsilon<R/10$ and $\log(R/\epsilon)\asymp L$. Localization
\eqref{eq:localization} and Lemma~\ref{lem:thin} give
\begin{align*}
 F(1+rz,Y)
 &\ge c L\sum_Q R^m h^{2m+2}(h^{-5m/2}M_Q)^q\\
 &=c L\sum_Q R^m h^{-(m+1)/2}M_Q^q\\
 &\ge c r^{-1}L\int_{\R^m}M(y,z)^q\dd y.
\end{align*}
Here $R^m\asymp r^m\delta^m$, $h=r^2$, and
$\int_QM^q\le\delta^mM_Q^q$.
Integrating over $t=1+rz$ and using \eqref{eq:totalupper},
\[
 c L\int_0^1\int_{\R^m}M^q
 \le\int_1^{1+r}\operatorname{avg}_{|Y|=h}F(t,Y)\dd t
 \le C E L^2.
\]
One logarithm cancels, proving Proposition~\ref{thm:main}.

\section{Arbitrary tube locations and the maximal operator}

First \eqref{eq:overlap} implies the global weighted tube estimate
\begin{equation}\label{eq:globaloverlap}
 \left\|\sum_j\lambda_j1_{T_\delta(\omega_j,a_j)}
             \right\|_{L^q(\R^n)}
 \le C_n L^{1/q}
               \left(\delta^m\sum_j\lambda_j^q\right)^{1/q}
\end{equation}
for arbitrary centers $a_j$ and $\delta$-separated
$\omega_j\in S^{n-1}$.
Here are the elementary localization details.
Partition the sphere into a fixed number of sufficiently small caps,
and rotate one cap to the vertical direction.
The corresponding slopes have size at most $\gamma$ and mutual
separation at least $c_n\delta$.
Partition the centers into unit lattice cubes.
In each cube, translate the cube to the origin. Every tube in that
group is contained in
\[
 \{(y,z): |y-b_j-zv_j|\le C_n\delta,\ z\in I\},
\]
where $I$ is a common interval of fixed length and $|b_j|\le C_n$.
A fixed dilation and a translation in $z$ put $I$ into $[0,1]$.
The thickening caused by this covering changes $\delta$ only by a
fixed factor. Partition the directions a further fixed number of
times so that their separations are at least that new radius.
Such a partition exists because a ball of fixed multiple of $\delta$
contains only a bounded number of the original separated directions.
Apply \eqref{eq:overlap} to each family. Each tube lies in a fixed
enlargement of its center cube; these enlarged cubes have bounded
overlap. The inequality
$(\sum_{\ell=1}^N u_\ell)^q\le N^{q-1}\sum_{\ell=1}^N u_\ell^q$
for bounded $N$ therefore permits summing the estimates. The same
inequality handles the finitely many direction partitions.
All dilation factors are dimensional constants, and
$\log(2/(c\delta))\asymp L$. This proves
\eqref{eq:globaloverlap} for small $\delta$.
For $\delta$ bounded below, the number of separated directions is
bounded, so the estimate is immediate.
Raising \eqref{eq:globaloverlap} to the power $q$ proves
Theorem~\ref{thm:tubes}.

We next prove Corollary~\ref{cor:maximal}.
Choose a maximal $\delta$-separated set $\{\omega_j\}$ on the sphere.
Assign each $\omega$ to a nearest $\omega_j$. Each resulting piece
has surface measure at most $C_n\delta^m$, and
\[
 T_\delta(\omega,a)\subset T_{C_n\delta}(\omega_j,a).
\]
Consequently
\[
 \|K_\delta f\|_{L^n(S^{n-1})}
 \le C_n\left(\delta^m\sum_j
                  K_{C_n\delta}f(\omega_j)^n\right)^{1/n}.
\]
The same finite direction partition used above allows
\eqref{eq:globaloverlap} for radius $C_n\delta$.
Test the discrete norm by coefficients $\lambda_j\ge0$ with
$\delta^m\sum_j\lambda_j^q=1$.
For any choices of centers, writing
$T_j=T_{C_n\delta}(\omega_j,a_j)$, H\"older gives
\[
 \delta^m\sum_j\frac{\lambda_j}{|T_j|}
                  \int_{T_j}|f|
 \le C_n\int_{\R^n}|f|\sum_j\lambda_j1_{T_j}
 \le C_nL^{1/q}\|f\|_{L^n(\R^n)}.
\]
There are finitely many directions, so centers may be chosen
independently to approach every supremum. Taking these suprema and
then the discrete dual supremum proves \eqref{eq:maximal}, since
$1/q=(n-1)/n$. Tube averages depend continuously on direction and
center for $f\in L^n$ by the $L^q$ continuity of the tube indicators;
their suprema are measurable.

\subsection*{The converse implication}

Suppose \eqref{eq:maximal} holds. Let $\omega_j$ be
$\delta$-separated, and choose pairwise disjoint spherical caps
$U_j$ of radius $c\delta$ around them, for a small dimensional $c$.
Then $\sigma(U_j)\asymp\delta^m$.
For each fixed $a_j$ and $\omega\in U_j$,
\[
 T_\delta(\omega_j,a_j)\subset T_{2\delta}(\omega,a_j).
\]
The ratio of tube volumes is bounded. Therefore, for every
$f\in L^n(\R^n)$,
\[
 \delta^m\sum_j\left(
       \frac1{|T_\delta(\omega_j,a_j)|}
           \int_{T_\delta(\omega_j,a_j)}|f|\right)^n
 \le C_n\|K_{2\delta}f\|_n^n
 \le C_nL^{n-1}\|f\|_n^n.
\]
For sufficiently small $\delta$ this uses the assumed maximal
estimate at $2\delta$; the remaining scales follow directly from
H\"older's inequality on a tube.
Discrete H\"older, with exponents $q,n$, now gives
\[
 \int |f|\sum_j\lambda_j1_{T_\delta(\omega_j,a_j)}
 \le C_n
 \left(\delta^m\sum_j\lambda_j^q\right)^{1/q}
       L^{1/q}\|f\|_n.
\]
The fact that all these tubes have the same volume, comparable to
$\delta^m$, is used in this step.
Duality in $L^q(\R^n)$ proves \eqref{eq:globaloverlap} and hence
\eqref{eq:maintubes}. This establishes the asserted equivalence.

\section{Volume and sharpness}

Specialize Theorem~\ref{thm:tubes} to $\lambda_j=1$.
Put $M=\sum_j1_{T_j}$, $V=\bigcup_jT_j$, and
$E=\delta^m\#\{j\}$. Then $\int M\asymp E$, so H\"older gives
\[
 E\lesssim\int M
 \le |V|^{1/(m+1)}\left(\int M^q\right)^{m/(m+1)}
 \le C_n|V|^{1/(m+1)}(EL)^{m/(m+1)}.
\]
Thus $|V|\ge c_nEL^{-m}$; the empty family is immediate.

For a Kakeya set $K$, choose a maximal $\delta$-separated set of
directions, of cardinality comparable to $\delta^{-m}$, and choose
one unit segment in $K$ in each of these finitely many directions.
Their $\delta$-neighborhoods are tubes $T_j\subset N_\delta(K)$.
Since $E\asymp1$, the union estimate gives
\eqref{eq:neighborhood}. No boundedness or measurability of $K$
is needed: $N_\delta(K)$ is open, and the inequality is automatic
if it has infinite volume.

\subsection*{Sharpness}

For the maximal logarithm, take $\delta>0$ sufficiently small and let
\[
 f_\delta(X)=|X|^{-1}1_{\{2\delta<|X|<1/2\}}.
\]
Then $\|f_\delta\|_n\asymp L^{1/n}$. Every unit tube centered at the
origin has average at least
$c_n\int_{4\delta}^{1/4}t^{-1}\dd t\gtrsim L$.
Thus $\|K_\delta f_\delta\|_n/\|f_\delta\|_n
\gtrsim L^{(n-1)/n}$.
For overlap sharpness, take a full separated grid of slopes in a
fixed cone and let all tubes start at the origin, with
$\lambda_T=1$. At height $z$ between $C\delta$ and a fixed small
constant, the multiplicity on an interior cone is comparable to
$z^{-m}$ and the time-slice volume is comparable to $z^m$.
Hence $\int M^q\gtrsim\int_{C\delta}^c z^{-1}\dd z\gtrsim L$
while $E\asymp1$.
For neighborhood sharpness, \cite[Theorem 1.3]{Construction}
constructs Kakeya sets with
$|N_\delta(K)|\lesssim L^{-(n-1)}$.

\section{Dependencies and scope}

The only nonlinear existence theorem invoked is the smooth Dirichlet
result of Ma--Zhu. Proposition~\ref{prop:dirichlet} checks its domain,
boundary data, and subsolution hypotheses on
$B^{3m}\times\mathbb T^m$ and explains smooth dependence on the
auxiliary parameters. The cofactor and Green identities, scalar
superharmonicity, source regularization, localization, weighted trace
bound, compact packets, scalar propagation, and thin-source estimate
are all supplied above.

The proof treats arbitrary finite families and arbitrary nonnegative
coefficients. Constants may depend on dimension; they are independent
of coefficient ratios, tube locations, and the number of tubes in
Theorem~\ref{thm:tubes}. The normalization in
Proposition~\ref{thm:main} temporarily introduces a bounded-intercept
constant, which is made uniform by the spatial decomposition in the
proof of \eqref{eq:globaloverlap}.

The complete mathematical deduction remains subject to independent
verification. The computational checks accompanying its development
test algebraic identities and scaling; they are not used as proof
of any lemma.

\clearpage
\begin{thebibliography}{9}
\bibitem{MZ}
X.-N. Ma and J. Zhu,
\emph{The Dirichlet problem for quaternionic Monge--Amp\`ere equations
on flat quaternionic quotients}, arXiv:2611.03214, November 2, 2026.
Theorem 1.1 and Sections 5--6.
\url{https://arxiv.org/abs/2611.03214}.
\bibitem{OpenAI}
OpenAI, \emph{Local smoothing for the Euclidean wave equation},
October 25, 2026. Source sections \emph{calculus.tex},
\emph{current.tex}, \emph{energy.tex}, and \emph{dirichlet.tex}.
\href{https://github.com/openai/math/tree/main/preprints/Local-smoothing-for-the-Euclidean-wave-equation-October-25-2026}{Manuscript and source}.
\bibitem{Draft}
M. Lewko, \emph{A quadratic positive-current argument for Minkowski
dimension in three dimensions}, project research draft,
November 10, 2026.
\bibitem{Addendum}
\emph{Recovering the sharp logarithm in the three-dimensional tube
argument}, project research addendum, November 11, 2026.
\bibitem{AllDimensions}
\emph{A sharp logarithmic tube estimate in every dimension},
project research draft, November 11, 2026.
\bibitem{KN}
S. Ko{\l}odziej and N. C. Nguyen,
\emph{The Dirichlet problem for the Monge--Amp\`ere equation on
Hermitian manifolds with boundary},
Calc. Var. Partial Differential Equations \textbf{62} (2023), article 1.
\url{https://doi.org/10.1007/s00526-022-02336-y}.
\bibitem{Construction}
R. J. Fern\'andez-Delgado and M. de la Salle,
\emph{A construction of Kakeya Sets in Arbitrary Dimension},
arXiv:2607.14824, version 3, September 17, 2026.
\url{https://arxiv.org/abs/2607.14824}.
\end{thebibliography}
\end{document}
